\problema{Max Area of Island}{695}{src/01-grafos/695-max-area-of-island.cpp}{M}

Nos dan una cuadrícula de \texttt{0} (agua) y \texttt{1} (tierra). Una
\emph{isla} es un grupo de tierras conectadas en horizontal o vertical, y su
\emph{área} es cuántas celdas de tierra la forman. Queremos el área de la
isla más grande (0 si no hay tierra).

\paragraph{Intuición.} Es el mismo sobrevuelo de \emph{Number of Islands},
pero ahora no contamos islas: medimos cada una. Al pisar tierra nueva
recorremos toda la isla conectada sumando 1 por celda y hundiéndola para no
recontarla. Nos quedamos con el área máxima vista.

\begin{center}
\ttfamily
\begin{tabular}{ccccc}
0 & 0 & 1 & 0 & 0\\
0 & 0 & 1 & 1 & 0\\
0 & 1 & 1 & 0 & 0\\
0 & 0 & 0 & 0 & 1\\
\end{tabular}
\end{center}

\noindent La isla central conecta \textbf{5} celdas; la celda suelta de abajo
a la derecha es otra isla de área 1. La respuesta es \textbf{5}.

\paragraph{Solución.} DFS que, en lugar de un valor booleano, devuelve el
tamaño de la isla: \texttt{1} por la celda actual más la suma recursiva de sus
cuatro vecinas. Hundimos cada celda visitada marcándola como agua.

\lstinputlisting{src/01-grafos/695-max-area-of-island.cpp}

\complejidad{$O(m\cdot n)$}{$O(m\cdot n)$}

\paragraph{Notas de entrevista (Amazon).} Si te acaban de preguntar
\emph{Number of Islands}, este es el seguimiento natural: cambia el DFS de
``contar'' a ``medir''. Comenta que hundir la cuadrícula ahorra el arreglo
\texttt{visited}; si no puedes mutar la entrada, usa BFS con cola o un
conjunto de visitados. El máximo se puede llevar como acumulador global.
